% =============================================================================
% Clifford Seminar Group — Thesenpapier Template
%
% How to use this template:
%   1. Copy this file to:
%        tex/thesenpapier/sessionNN-<slug>/sessionNN-thesenpapier.tex
%   2. Replace the metadata fields (\seminartitle, \seminarsession, etc.)
%      and section contents below.
%   3. Compile with:
%        latexmk -pdf sessionNN-thesenpapier.tex
%      (or: pdflatex -> biber -> pdflatex -> pdflatex)
%
% Relative paths to clifford-seminar.sty, assets/logo.jpg, and
% bibliography/clifford-seminar-group.bib work identically here and in any
% tex/thesenpapier/sessionNN-<slug>/ directory.
% =============================================================================
\documentclass[11pt,a4paper]{article}

\usepackage[thesenpapier]{../../templates/clifford-seminar}

% -----------------------------------------------------------------------------
% Document Metadata (Edit for each session)
% -----------------------------------------------------------------------------
\seminartitle[Universal Clifford Construction]{%
  Universal Clifford Construction, Involutions, and the Geometric Product%
}
\seminarsubtitle{Template Position Paper --- Replace with Session Subtitle}
\seminarsession{Session NN}
\seminarmodule{Module II: Continuous Symmetries and Clifford Algebras}
\seminardate{DD Month YYYY}
\seminarinstitution{Universidad Distrital Francisco Jos\'e de Caldas}

\seminarauthors{Lead Author (Senior) \and Co-Author (Junior)}
\seminarreferat{Senior Name \& Junior Name}
\seminarkorreferat{Respondent Name}

\begin{document}

\makeseminarheader

\begin{seminarabstract}[Conceptual Scope]
  Summarize the physical motivation and mathematical objective of this
  session in three to five concise sentences. Explain why the algebraic
  structure introduced here is required by the guiding physical problem rather
  than introduced in isolation, and indicate what participants are expected
  to prepare prior to the meeting.
\end{seminarabstract}

% =============================================================================
\section{Physical Motivation and Conceptual Context}
\label{sec:motivation}
% =============================================================================

Open the position paper by grounding the algebraic machinery in a concrete
physical problem (for example, orbital motion in central-force dynamics,
measurement duality, or relativistic kinematics). Contrast the invariant
geometric formulation with standard coordinate- or matrix-based treatments
\cite{hestenesOerstedMedalLecture2003,doranGeometricAlgebraPhysicists2013}.

In a bare real linear space $V$, elements can be added and scaled, yet no
intrinsic associative product exists between vectors that simultaneously
encodes metric projection and oriented planar area. Equipping $V$ with a
non-degenerate quadratic form $q : V \to \R$ supplies the metric input needed
to pass from the free tensor algebra $T(V)$ to the universal Clifford algebra
$\Cl(V, q)$ \cite[Chap.~2]{vazIntroductionCliffordAlgebras2016}.

% =============================================================================
\section{Foundational Definitions and Algebraic Structure}
\label{sec:definitions}
% =============================================================================

State the primary axioms, definitions, and structural propositions that will
anchor the 40-minute blackboard exposition. Keep notation consistent with the
shared macros provided by the style file (e.g., $\Cl(V,q)$,
$\pair{\cdot}{\cdot}$, $\rev{u}$, $\ginv{u}$, $\clconj{u}$).

\begin{definition}[Universal Clifford Algebra]
  \label{def:clifford-algebra}
  Let $(V, q)$ be a finite-dimensional real quadratic space, and let
  $T(V) = \bigoplus_{r=0}^{\infty} T^r(V)$ denote its free tensor algebra.
  Let $\mathcal{I}_q \subset T(V)$ be the two-sided ideal generated by elements
  of the form $v \otimes v - q(v) 1$ for all $v \in V$. The
  \emph{Clifford algebra} of $(V, q)$ is the quotient associative unital
  algebra
  \begin{equation}
    \label{eq:clifford-quotient}
    \Cl(V, q) = T(V) / \mathcal{I}_q.
  \end{equation}
\end{definition}

\begin{proposition}[Fundamental Contraction and Polarization]
  \label{prop:polarization}
  For all vectors $u, v \in V \subset \Cl(V, q)$, the Clifford (geometric)
  product satisfies the anticommutation identity
  \begin{equation}
    \label{eq:fundamental-identity}
    uv + vu = 2\, b(u, v) 1,
  \end{equation}
  where $b(u, v) = \tfrac{1}{2}\bigl(q(u + v) - q(u) - q(v)\bigr)$ is the
  symmetric bilinear form associated with $q$.
\end{proposition}

\begin{proof}
  Apply the defining relation $w^2 = q(w) 1$ in $\Cl(V, q)$ to the sum
  $w = u + v$:
  \[
    q(u + v) 1 = (u + v)^2 = u^2 + uv + vu + v^2 = \bigl(q(u) + q(v)\bigr) 1 + uv + vu.
  \]
  Rearranging and invoking the polarization identity yields
  \cref{eq:fundamental-identity}.
\end{proof}

\begin{keyresult}[Core Decomposition of the Geometric Product]
  For any vectors $u, v \in V \subset \Cl(V, q)$, the associative geometric
  product decomposes uniquely into a symmetric scalar part and an
  antisymmetric bivector part:
  \begin{equation}
    \label{eq:geometric-product}
    uv = u \cdot v + u \wedge v, \qquad
    u \cdot v = \frac{1}{2}(uv + vu), \quad
    u \wedge v = \frac{1}{2}(uv - vu).
  \end{equation}
\end{keyresult}

\subsection{Involutions and Grading Summary}
\label{subsec:involutions}

Use clean \texttt{booktabs} tables (see \cref{tab:involutions}) to summarize
structural identities, sign conventions, or correspondences for rapid
reference during the session.

\begin{table}[htbp]
  \centering
  \caption{Principal automorphisms and anti-automorphisms on a homogeneous
    $k$-blade $A_k \in \Cl^k(V, q)$.}
  \label{tab:involutions}
  \small
  \renewcommand{\arraystretch}{1.22}
  \begin{tabularx}{\textwidth}{@{} >{\bfseries}l c c >{\raggedright\arraybackslash}X @{}}
    \toprule
    \textbf{Operation} & \textbf{Notation} & \textbf{Action on $A_k$} & \textbf{Algebraic Character} \\
    \midrule
    Grade involution & $\ginv{A}_k$ & $(-1)^k A_k$ &
      Algebra automorphism; separates even ($\Cl^+$) and odd ($\Cl^-$) parts \\
    Reversion & $\rev{A}_k$ & $(-1)^{k(k-1)/2} A_k$ &
      Anti-automorphism reversing the order of vector factors \\
    Clifford conjugation & $\clconj{A}_k$ & $(-1)^{k(k+1)/2} A_k$ &
      Composition $\clconj{A} = \ginv{\rev{A}}$; anti-automorphism \\
    \bottomrule
  \end{tabularx}
\end{table}

\begin{remark}
  While $\Cl(V, q)$ is only $\Z_2$-graded as an algebra (because $v^2 = q(v)1$
  mixes degrees differing by $2$), it remains canonically isomorphic as a
  $\Z$-graded linear space to the exterior algebra $\bigwedge V$.
\end{remark}

% =============================================================================
\section{Derivation Agenda for the Blackboard Exposition}
\label{sec:agenda}
% =============================================================================

Outline the sequential derivation stages planned for the 40-minute
\emph{Referat} so participants can track the logical arc in advance:
\begin{enumerate}
  \item \textbf{Stage 1 (10 min):} Universal property of $T(V)$ and quotient
    construction of $\Cl(V, q) = T(V)/\mathcal{I}_q$.
  \item \textbf{Stage 2 (10 min):} Proof of the universal property of
    $\Cl(V, q)$ for Clifford maps $f : V \to A$ satisfying $f(v)^2 = q(v)1_A$.
  \item \textbf{Stage 3 (10 min):} Construction of grade involution,
    reversion, and Clifford conjugation directly from the universal property.
  \item \textbf{Stage 4 (10 min):} Decomposition of the geometric product
    $uv = u \cdot v + u \wedge v$ and explicit basis of $\Cl(3,0)$.
\end{enumerate}

% =============================================================================
\section{Theses for Plenary Discussion}
\label{sec:theses}
% =============================================================================

Formulate three to four crisp, falsifiable or debate-provoking theses for the
\emph{Korreferat} and plenary discussion. Each thesis should take a clear
conceptual stance on mathematical structure or physical interpretation.

\begin{thesis}[Primacy of the Geometric Product]
  The dot product and cross product of elementary vector analysis are not
  independent primitive operations, but coordinate-dependent projections of a
  single associative geometric product $uv = u \cdot v + u \wedge v$ that is
  invertible for all non-null vectors.
\end{thesis}

\begin{thesis}[Coordinate-Free Involutions]
  Grade involution, reversion, and Clifford conjugation are intrinsic
  structural morphisms induced solely by the universal property of $\Cl(V,q)$,
  requiring neither matrix representations nor an orthonormal basis for their
  definition.
\end{thesis}

\begin{thesis}[Physical Meaning of Oriented Blades]
  Replacing axial vectors (such as angular momentum $\vec{L} = \vec{r} \times
  \vec{p}$) with intrinsic bivectors ($L = r \wedge p$) eliminates artificial
  three-dimensional artifacts and extends seamlessly to relativistic
  spacetime $\Cl(1,3)$.
\end{thesis}

% =============================================================================
\section{Preparatory Exercises}
\label{sec:exercises}
% =============================================================================

Participants are expected to attempt the following exercises prior to the
session and submit at least one written question 24~hours before the meeting.

\begin{enumerate}
  \item Verify directly from \cref{def:clifford-algebra} that two vectors
    $u, v \in V$ are orthogonal with respect to $b$ if and only if they
    anticommute in $\Cl(V, q)$, i.e., $uv = -vu$.
  \item Let $\{e_1, e_2, e_3\}$ be an orthonormal basis of $\R^{3,0}$ with
    $e_i^2 = +1$.
    \begin{enumerate}
      \item Compute the square of the unit pseudoscalar $I_3 = e_1 e_2 e_3$
        and show that $I_3$ commutes with every element of $\Cl(3,0)$.
      \item Demonstrate that the even subalgebra $\Cl^+(3,0)$ is isomorphic
        to Hamilton's quaternion algebra $\mathbb{H}$.
    \end{enumerate}
  \item Use the universal property of $\Cl(V, q)$ to prove the existence and
    uniqueness of the grade involution $\ginv{u}$ without referring to any
    basis of $V$.
\end{enumerate}

% =============================================================================
\section{Preparatory Literature}
\label{sec:literature}
% =============================================================================

Participants should study the following primary references ahead of the
session:
\begin{itemize}
  \item Quotient construction, universal property, and involutions:
    Vaz and da Rocha \cite[Chap.~2]{vazIntroductionCliffordAlgebras2016}.
  \item Geometric product and physical interpretation in three dimensions:
    Doran and Lasenby \cite[Chaps.~1--2]{doranGeometricAlgebraPhysicists2013};
    Hestenes \cite{hestenesNewFoundationsClassical2002}.
\end{itemize}

\printbibliography[title={References}]

\end{document}
